 %%%%%%%%%%%%%%%%%%%%%%%%%%%%%%%%%%%%%%%%%
% Beamer Presentation
% LaTeX Template
% Version 1.0 (10/11/12)
%
% This template has been downloaded from:
% http://www.LaTeXTemplates.com
%
% License:
% CC BY-NC-SA 3.0 (http://creativecommons.org/licenses/by-nc-sa/3.0/)
%
%%%%%%%%%%%%%%%%%%%%%%%%%%%%%%%%%%%%%%%%%

%----------------------------------------------------------------------------------------
%	PACKAGES AND THEMES
%----------------------------------------------------------------------------------------

\documentclass[9pt,aspectratio=169]{beamer}
\usepackage[utf8]{inputenc}
\usepackage{bm}
\usepackage{bbm}
\usepackage{graphicx}
\usepackage{cmbright}
\usefonttheme{professionalfonts}
\DeclareSymbolFont{greekletters}{OML}{sansmath}{m}{it}
\DeclareMathSymbol{\beta}{\mathord}{greekletters}{"0C}
%\usefonttheme{serif}
% The Beamer class comes with a number of default slide themes
% which change the colors and layouts of slides. Below this is a list
% of all the themes, uncomment each in turn to see what they look like.

%\usetheme{default}
%\usetheme{AnnArbor}
% \usetheme{Antibes}
%\usetheme{Bergen}
%\usetheme{Berkeley}
%\usetheme{Berlin}
%\usetheme{Boadilla}
% \usetheme{CambridgeUS}
%\usetheme{Copenhagen}
%\usetheme{Darmstadt}
% \usetheme{Dresden}
\usetheme{Frankfurt}
%\usetheme{Goettingen}
% \usetheme{Hannover}
% \usetheme{Ilmenau}
% \usetheme{JuanLesPins}
%\usetheme{Luebeck}
%\usetheme{Madrid}
%\usetheme{Malmoe}
%\usetheme{Marburg}
% \usetheme{Montpellier}
% \usetheme[width=2.5cm]{PaloAlto}
%\usetheme{Pittsburgh}
% \usetheme{Rochester}
%\usetheme{Singapore}
% \usetheme{Szeged}
% \usetheme{Warsaw}
% \setbeamertemplate{title page}[default][colsep=-4bp,rounded=false]
\setbeamertemplate{headline}{}
\defbeamertemplate{footline}{myframe number}
{
  \hspace{1pt}
    \usebeamercolor[structure]{page number in head/foot}%
  \usebeamerfont{page number in head/foot}%
  \raisebox{1.7ex}[0pt][0pt]{% <--- change here
    \insertframenumber\,/\,\inserttotalframenumber\kern1em}%
}
\setbeamertemplate{footline}[myframe number]
\AtBeginSection[]
{
  \begin{frame}
    \frametitle{Outline}
    \begin{minipage}{1.0\linewidth}
      \tableofcontents[currentsection,currentsubsection]
    \end{minipage}
  \end{frame}
}

\AtBeginSubsection[]
{
  \begin{frame}
    \frametitle{Outline}
    \begin{minipage}{1.0\linewidth}
      \tableofcontents[currentsection,currentsubsection]
    \end{minipage}
  \end{frame}
}

% As well as themes, the Beamer class has a number of color themes
% for any slide theme. Uncomment each of these in turn to see how it
% changes the colors of your current slide theme.
%\definecolor{mycolor}{rgb}{.8,.0,.3}
%\setbeamercolor*{palette primary}{use=structure,fg=white,bg=mycolor}
%\usecolortheme{albatross}
%\usecolortheme{beaver}
%\usecolortheme{beetle}
%\usecolortheme{crane}
%\usecolortheme{dolphin}
%\usecolortheme{dove}
%\usecolortheme{fly}
%\usecolortheme{lily}
%\usecolortheme{orchid}
% \usecolortheme{rose}
%\usecolortheme{seagull}
\usecolortheme{seahorse}
%\usecolortheme{whale}
%\usecolortheme{wolverine}
%\usecolortheme[named=mycolor]{structure}
%\setbeamertemplate{footline} % To remove the footer line in all slides uncomment this line
%\setbeamertemplate{footline}[page number] % To replace the footer line in all slides with a simple slide count uncomment this line

%\setbeamertemplate{navigation symbols}{} % To remove the navigation symbols from the bottom of all slides uncomment this line


\usepackage{booktabs} % Allows the use of \toprule, \midrule and \bottomrule in tables

% Custom commands
\newcommand{\R}{\mathbb{R}}
\newcommand{\Z}{\mathbb{Z}}

\title{Lecture 8: Infinite-Horizon Optimization and Dynamic Programming}
\date{}

\author{Junnan Zhang} % Your name
\institute[Paula and Gregory Chow Institute for Studies in Economics\\
Xiamen University] % Your institution as it will appear on the bottom of every slide, may be shorthand to save space
{
	Paula and Gregory Chow Institute for Studies in Economics \\ Xiamen University  \\ % Your institution for the title page
	\medskip
}
\date{Fall, 2026} % Date, can be changed to a custom date
\setbeamertemplate{itemize items}[default]
\setbeamertemplate{sidebar}[sidebar right]
% \definecolor{myco}{HTML}{a67678}
\definecolor{myco}{HTML}{8776a6}
\setbeamercolor{itemize item}{fg=myco} % all frames will have red bullets
\setbeamercolor{itemize subitem}{fg=myco} % all frames will have red bullets
\setbeamercolor{block title}{bg=myco}
\setbeamercolor{structure}{fg=myco}
% \setbeamercovered{transparent}


\begin{document}

\frame{\titlepage}

\begin{frame}{Chapter Overview}
    \begin{itemize}
        \item Brief introduction to infinite-horizon optimization in discrete
        time under certainty
        \item Main purpose: introduce the theory and techniques of dynamic
        programming
    \end{itemize}
\end{frame}

\begin{frame}{Chapter Structure}
The material is presented in four parts:
\begin{enumerate}
    \item Introduces the problem and provides theoretical results necessary for
    applications of stationary dynamic programming techniques
    \item Additional mathematical tools
    \item Applications
\end{enumerate}
\end{frame}

\section{Problem Description}

\begin{frame}{Discrete-Time Infinite-Horizon Optimization}
The canonical discrete-time infinite-horizon optimization program:
\begin{align*}
&\sup_{\{x_t,y_t\}_{t=0}^{\infty}} \sum_{t=0}^{\infty} \beta^t \tilde{U}(t, x_t, y_t)\\
\text{s.t. }
&y_t \in \tilde{G}(t, x_t) \text{ for all } t \geq 0\\
&x_{t+1} = \tilde{f}(t, x_t, y_t) \text{ for all } t \geq 0\\
&x_0 \text{ given}
\end{align*}
\small
\begin{itemize}
    \item $\beta \in [0,1)$ is the discount factor
    \item $x_t \in X \subset \R^{K_x}$ and $y_t \in Y \subset \R^{K_y}$ for some $K_x, K_y \geq 1$
    \item $x_t$: state variables (state vector)
    \item $y_t$: control variables (control vector)
    \item $\tilde{U}: \Z_+ \times X \times Y \to \R$ is the instantaneous payoff function
    \item $\tilde{G}(t,x)$ is a set-valued mapping (correspondence): $\tilde{G}: \Z_+ \times X \rightrightarrows Y$
\end{itemize}
\end{frame}

\begin{frame}{Constraints and Evolution}
    \begin{itemize}
        \item First constraint: $y_t \in \tilde{G}(t, x_t)$ specifies what
        values of control vector $y_t$ are allowed, given state $x_t$ at time
        $t$
        \item Evolution equation: $\tilde{f}: \Z_+ \times X \times Y \to X$
        specifies evolution of state vector as function of last period's
        state and control vectors
        \item This formulation highlights distinction between state and control variables
        \item Often more convenient to work with transformation to eliminate $y_t$
    \end{itemize}
\end{frame}

\begin{frame}{Problem 6.1: Transformed Formulation}
    \begin{align*}
        V^*(0, x_0) &= \sup_{\{x_t\}_{t=0}^{\infty}}
        \sum_{t=0}^{\infty} \beta^t U(t, x_t, x_{t+1})\\ 
        \text{s.t. } &x_{t+1} \in G(t, x_t) \text{ for all } t \geq 0\\
        &x_0 \text{ given}
    \end{align*}
    \begin{itemize}
        \item $x_t \in X \subset \R^K$ (now $K = K_x$)
        \item $x_t$ corresponds to state vector, $x_{t+1}$ plays role of control vector at time $t$
        \item $U: \Z_+ \times X \times X \to \R$ is instantaneous payoff function
        \item $G: \Z_+ \times X \rightrightarrows X$ specifies constraint correspondence
    \end{itemize}
\end{frame}

\begin{frame}{Value Function}
\begin{itemize}
    \item $V^*: \Z_+ \times X \to \R$ is the \textbf{value function}, which
    specifies supremum (highest possible value) that objective function can
    reach starting with some $x_t$ at time $t$
    \item When maximal value is attained by sequence
    $\{x^*_{t+1}\}_{t=0}^{\infty} \in X^{\infty}$, we refer to this as a
    \textbf{solution} or an \textbf{optimal plan}
    \item Objective: characterize optimal plan $\{x^*_{t+1}\}_{t=0}^{\infty}$
    and value function $V^*(0, x_0)$
    \item Here, $X^{\infty}$ is countable infinite product of set $X$. If
    $X$ is bounded, then $X^{\infty} \subset \ell^{\infty}$, where
    $\ell^{\infty}$ is the vector space of infinite sequences that are bounded
    with the sup norm $\|\cdot\|_{\infty}$ (or simply $\|\cdot\|$)
\end{itemize}
\end{frame}

% \begin{frame}{More General Formulation}
% \begin{itemize}
% \item More general formulation would involve undiscounted objective:
% $$\sup_{\{x_t\}_{t=0}^{\infty}} U_{x(0}, x_{1}, \ldots)$$
% \item This added generality not particularly useful for most economic growth problems
% \item Discounted time-separable objective function ensures \textbf{time consistency} of optimal plan
% \item Time consistency discussed in previous chapter (Exercise 5.1)
% \item Problem 6.1 somewhat abstract but has advantage of being tractable
% \item General enough to nest many interesting economic applications
% \end{itemize}
% \end{frame}

\begin{frame}{Example 6.1: Optimal Growth Problem}
Consider an optimal growth problem:
\begin{align*}
&\max_{\{k_t,c_t\}_{t=0}^{\infty}} \sum_{t=0}^{\infty} \beta^t u(c_t)\\
\text{s.t. } &k_{t+1} = f(k_t) + (1-\delta)k_t - c_t\\
&k_t \geq 0, \text{ given } k_0 > 0, \text{ and } u: \R_+ \to \R
\end{align*}
Let $x_t = k_t$ and $x_{t+1} = k_{t+1}$. The objective function can be
written as:
\begin{equation*}
\max_{\{k_t\}_{t=0}^{\infty}} \sum_{t=0}^{\infty} \beta^t u(f(k_t) - k_{t+1} + (1-\delta)k_t)
\end{equation*}
subject to $k_t \geq 0$. This is special case of Problem 6.1 with:
\begin{itemize}
\item $U(t, k_t, k_{t+1}) = u(f(k_t) - k_{t+1} + (1-\delta)k_t)$
\item Constraint correspondence $G(t, k_t) = [0, f(k_t) + (1-\delta)k_t]$ (since $c_t \geq 0$)
\end{itemize}
\end{frame}

\begin{frame}{Stationarity Feature}
\begin{itemize}
    \item Notable feature emphasized by Example 6.1: $U$ and $G$ do not
    explicitly depend on time
    \item This feature is fairly common in economics and many interesting
    problems can be formulated in \textbf{stationary form}
    \item Stationary problem involves:
    \begin{itemize}
        \item Objective function that is a discounted sum
        \item $U$ and $G$ functions that do not explicitly depend on time
    \end{itemize}
\end{itemize}
\end{frame}

\section{Stationary Dynamic Programming}

\begin{frame}{Stationary Dynamic Programming}
    Consider the stationary form of Problem 6.1.
    \vspace{0.5cm}
    
    \textbf{Problem 6.2}:
    \begin{align*}
        V^*(x_0) &= \sup_{\{x_t\}_{t=0}^{\infty}} \sum_{t=0}^{\infty} \beta^t U(x_t, x_{t+1})\\
        \text{s.t. } &x_{t+1} \in G(x_t) \text{ for all } t \geq 0\\
        &x_0 \text{ given}
    \end{align*}
\begin{itemize}
\item Constraint correspondence: $G: X \rightrightarrows X$
\item Instantaneous payoff functions: $U: X \times X \to \R$
\item Value function without time argument: $V^*(x_0)$
\end{itemize}
\end{frame}

\begin{frame}{Sequence vs. Functional Equation}
\begin{itemize}
    \item Problem 6.2 corresponds to the \textbf{sequence problem}: involves
    choosing infinite sequence $\{x_t\}_{t=0}^{\infty} \in X^{\infty}$
    \item Sequence problems sometimes have nice features, but solutions often
    difficult to characterize analytically and numerically
    \item Basic idea of dynamic programming: turn sequence problem into
    \textbf{functional equation}
    \item Transform problem into one of finding a function rather than a
    sequence
\end{itemize}
\end{frame}

\begin{frame}{The Bellman Equation}
\textbf{Problem 6.3}:
$$V(x) = \sup_{y \in G(x)} \{U(x,y) + \beta V(y)\}, \quad \text{for all } x \in X$$
where $V: X \to \R$.

\begin{itemize}
    \item This functional equation is also called the \textbf{``Bellman
      equation''} after Richard Bellman
    \item Instead of explicitly choosing a sequence
    $\{x_t\}_{t=0}^{\infty}$, we choose a \textbf{policy} $y \in G(x)
    \subset X$ for given $x \in X$
    \item Since $U(\cdot,\cdot)$ does not depend on time, no reason for
    policy to be time-dependent either
    \item Use $V$ for function defined in Problem 6.3 and $V^*$ for function
    defined in Problem 6.2
\end{itemize}
\end{frame}

\begin{frame}{Economic Interpretation}
$$V(x) = \sup_{y \in G(x)} \{U(x,y) + \beta V(y)\}, \quad \text{for all } x \in X$$
\begin{itemize}
    \item This is a \textbf{``recursive formulation''}: function $V(x)$
    appears on both left- and right-hand sides, defined recursively
    \item In many cases, solution to Problem 6.3 is simpler to characterize
    analytically than the corresponding solution to sequence problem
    \item Question: when are the two problems equivalent?
\end{itemize}
\end{frame}

\begin{frame}{Natural Derivation of the Bellman Equation}
    Suppose Problem 6.2 has maximum starting at $x_0$ attained by optimal
    sequence $\{x^*_t\}_{t=0}^{\infty}$ with $x^*_0 = x_0$. Under weak
    technical conditions:
    \begin{align*}
        V^*(x_0) &= \sum_{t=0}^{\infty} \beta^t U(x^*_t, x^*_{t+1})\\
        &= U(x_0, x^*_1) + \beta \sum_{s=0}^{\infty} \beta^s U(x^*_{s+1}, x^*_{s+2})\\
        &= U(x_0, x^*_1) + \beta V^*(x^*_1)
    \end{align*}
    This encapsulates the basic idea of dynamic programming:
    \textbf{Principle of Optimality}, which states that optimal plan can be
    broken into two parts:
    \begin{enumerate}
        \item What is optimal today
        \item The optimal continuation path
    \end{enumerate}
\end{frame}

\begin{frame}{Optimal Policy}
$$V(x) = \sup_{y \in G(x)} \{U(x,y) + \beta V(y)\}, \quad \text{for all } x \in X$$
\begin{itemize}
    \item Solution can be represented by a time invariant \textbf{policy
      function} (or policy mapping) $\pi: X \to X$ that determines which
    value of $x_{t+1}$ to choose for given state variable $x_t$
    \item Two complications:
    \begin{enumerate}
        \item Control reaching optimal value may not exist (reason for ``sup'' notation)
        \item Solution may involve policy \textbf{correspondence} $\Pi: X \rightrightarrows X$ (more than one maximizer)
    \end{enumerate}
    \item Ignore these complications for now
\end{itemize}
\end{frame}

\begin{frame}{Characterizing Policy Function}
    Once value function $V$ is determined, policy function is straightforward
    to characterize.

    If the optimal policy is given by $\pi(x)$, then:
    $$V(x) = U(x, \pi(x)) + \beta V(\pi(x)) \quad \text{for all } x \in X$$

    \begin{itemize}
        \item This equation follows from the fact that $\pi(x)$ is the
        optimal policy
        \item When $y = \pi(x)$, the right-hand side of the Bellman equation
        reaches maximal value $V(x)$
        \item This is one way of determining the policy function
    \end{itemize}
\end{frame}

\begin{frame}{Usefulness of Recursive Formulation}
\begin{itemize}
    \item There are powerful tools that establish existence of solution and
    some of its properties
    \item Next section: present results on relationship between solution to
    the sequence problem (Problem 6.2) and the recursive formulation (Problem
    6.3)
    \item Establish results on concavity, monotonicity, and differentiability
    of the value function
\end{itemize}
\end{frame}

\section{Stationary Dynamic Programming Results}

\begin{frame}{Stationary Dynamic Programming Theorems}
    Main purpose:
    \begin{itemize}
    \item Ensure sequence $\{x^*_t\}_{t=0}^{\infty} \in
    X^{\infty}$ that attains supremum in Problem 6.2 satisfies recursive
    equation of dynamic programming
    \begin{equation*}
        V(x^*_t) = U(x^*_t, x^*_{t+1}) + \beta V(x^*_{t+1}), \quad \text{for
          all } t = 0, 1, \ldots
    \end{equation*}
    \item The solution to the Bellman equation is also a solution to Problem 6.2
    \item In other words, establish equivalence results between solutions to
    Problems 6.2 and 6.3
\end{itemize}
\end{frame}

\begin{frame}{Feasible Sequences}
Define set of feasible sequences or plans starting with initial value $x_t$:
$$\Phi(x_t) = \{\{x_{s}\}_{s=t}^{\infty} : x_{s+1} \in G(x_s) \text{ for } s = t, t+1, \ldots\}$$

\begin{itemize}
    \item Intuitively, $\Phi(x_t)$ is the set of feasible choices of vectors
    starting from $x_t$
    \item Denote typical element of $\Phi(x_0)$ by $\mathbf{x} = (x_0,
    x_{1}, \ldots) \in \Phi(x_0)$
    \item Want sequence to satisfy:
    $$V(x^*_t) = U(x^*_t, x^*_{t+1}) + \beta V(x^*_{t+1}) \text{ for all } t = 0,1,\ldots$$
\end{itemize}
\end{frame}

\begin{frame}{Basic Assumptions}
\begin{block}{Assumption 6.1}
$G(x)$ is nonempty for all $x \in X$; and for all $x_0 \in X$ and $\mathbf{x} \in \Phi(x_0)$,
$$\lim_{n \to \infty} \sum_{t=0}^n \beta^t U(x_t, x_{t+1}) \text{ exists and is finite}$$
\end{block}

\begin{itemize}
    \item Assumption stronger than necessary for theory: for much of dynamic
    programming theory, it is sufficient that the limit exists
    \item But in economic applications, we are not interested in optimization
    problems where agents achieve infinite value
\end{itemize}
\end{frame}

\begin{frame}{Equivalence of Values}
    \begin{block}{Theorem 6.1 (Equivalence of Values)}
        Suppose Assumption 6.1 holds. Then for any $x \in X$, any solution
        $V^*(x)$ to Problem 6.2 is also a solution to Problem 6.3. Moreover,
        any solution $V(x)$ to Problem 6.3 is also a solution to Problem 6.2,
        so that $V^*(x) = V(x)$ for all $x \in X$.
    \end{block}

    \begin{itemize}
        \item Under Assumption 6.1, both sequence problem and recursive
        formulation achieve the same value
        \item Fundamental result that justifies the dynamic programming
        approach
    \end{itemize}
\end{frame}

\begin{frame}{Principle of Optimality}
\begin{block}{Theorem 6.2 (Principle of Optimality)}
Suppose Assumption 6.1 holds. Let $\mathbf{x}^* \in \Phi(x_0)$ be a feasible
plan that attains $V^*(x_0)$ in Problem 6.2. Then
\begin{equation*}
    \tag{6.3}
    V^*(x^*_t) = U(x^*_t, x^*_{t+1}) + \beta V^*(x^*_{t+1})    
\end{equation*}
for $t = 0, 1, \ldots$, with $x^*_0 = x_0$.

Moreover, if any $\mathbf{x}^* \in \Phi(x_0)$ satisfies equation (6.3), then
it attains the optimal value in Problem 6.2.
\end{block}
\begin{itemize}
  \item If any feasible plan satisfies the Bellman equation, then it attains
  optimal value
  \item We can go from solution of the Bellman equation to solution of the
  sequence problem and vice versa
\end{itemize}
\end{frame}

\begin{frame}{Example: Optimal Growth}
    \textbf{Example 6.4}: Consider an optimal growth model with log
    preferences, Cobb-Douglas technology, and full depreciation:
    \begin{align*}
        \max_{\{k_t,c_t\}_{t=0}^{\infty}} &\sum_{t=0}^{\infty} \beta^t \log c_t \\
        \text{s.t. } &k_{t+1} = k_t^{\alpha} - c_t \\
        &k_0 > 0
    \end{align*}
    where $\beta \in (0,1)$, $k$ is capital-labor ratio, and the resource
    constraint follows from the production function $K^{\alpha}L^{1-\alpha}$
    in per capita terms.
\end{frame}

\begin{frame}{Closed-Form Solution}
We solve the Bellman equation
\[
V(x) = \max_{y\ge 0} \{\log(x^\alpha - y) + \beta V(y)\}.
\]

Standard approach for log–Cobb–Douglas:
\begin{itemize}
    \item Conjecture that the value function is logarithmic:
    \[
    V(x) = A + B\log x.
    \]
    \item Use the Bellman equation to solve for unknown constants $(A,B)$.
    \item Obtain closed-form policy and value functions.
\end{itemize}
\end{frame}

\begin{frame}{Closed-Form Solution}
    \begin{itemize}
        \item Guess:
        \[ V(x)=A+B\log x.
        \]
        \item Bellman equation:
        \[
            A + B\log x=
            \max_{y}
            \left\{
                \log(x^\alpha - y) + \beta\left[A + B\log y \right]
            \right\}.
        \]
        \item FOC w.r.t.\ $y$:
        \[
            -\frac{1}{x^\alpha - y} + \frac{\beta B}{y} = 0
            \implies y = \frac{\beta B}{1+\beta B}x^\alpha.
        \]
        Using $B=\frac{\alpha}{1-\alpha\beta}$ gives
        $y=\alpha\beta x^\alpha$.
        \item Optimal next-period capital:
        \[
            k_{t+1} = y^*(x) = \alpha\beta x^\alpha.
        \]
    \end{itemize}
\end{frame}


\begin{frame}
    \frametitle{Closed-Form Solution}
    \begin{itemize}
        \item Substitute $y^*(x)=\alpha\beta x^\alpha$ into Bellman equation:
        \begin{align*}
            A+B\log x &= \log\left((1-\alpha\beta)x^\alpha\right) +
            \beta\left[ A + B\log(\alpha\beta x^\alpha) \right]\\
            &= \log(1-\alpha\beta) + \alpha\log x + \beta A + \beta
            B\left(\log(\alpha\beta)+\alpha\log x\right).
        \end{align*}
        \item Match coefficients on $\log x$:
        \[ B = \alpha + \beta B \alpha \implies
            B = \frac{\alpha}{1-\alpha\beta}
        \]
        \item Match constant terms:
        \[
            A = \log(1-\alpha\beta) + \beta A + \beta B\log(\alpha\beta).
        \]
        Solve:
        \[
            A=\frac{1}{1-\beta}\left[ \log(1-\alpha\beta) +
                \frac{\alpha\beta}{1-\alpha\beta}\log(\alpha\beta) \right].
        \]
    \end{itemize}
\end{frame}


\begin{frame}{Continuity and Compactness Assumptions}
\begin{block}{Assumption 6.2}
$X$ is a compact subset of $\R^K$, $G$ is nonempty-valued, compact-valued, and continuous. Moreover, $U: X_G \to \R$ is continuous, where $X_G = \{(x,y) \in X \times X : y \in G(x)\}$.
\end{block}

\begin{itemize}
    \item Most restrictive: the state space is compact
    \item Since $X$ is compact and $G(x)$ is compact-valued, $X_G$ is also compact
    \item Continuous function on compact domain: $U$ is bounded
\end{itemize}
\end{frame}


\begin{frame}{Existence of Solutions}
\begin{block}{Theorem 6.3 (Existence of Solutions)}
    Suppose that Assumptions 6.1 and 6.2 hold. Then there exists a unique
    continuous and bounded function $V: X \to \R$ that satisfies the Bellman
    equation. Moreover, for any $x_0 \in X$, an optimal plan $\mathbf{x}^*
    \in \Phi(x_0)$ exists.
\end{block}

Two major results:
\begin{enumerate}
    \item \textbf{Existence and uniqueness} of the value function
    \item \textbf{Existence} of an optimal plan
\end{enumerate}

Combined with Theorem 6.1: optimal policy function achieving supremum $V^*$ in Problem 6.2 exists and $V^*$ is continuous and bounded.
\end{frame}

\begin{frame}{Non-uniqueness of Optimal Plans}
\begin{itemize}
    \item Optimal plan in Problem 6.2 (or 6.3) may not be unique, even though
    the value function is unique: for example, when two alternative feasible
    sequences achieve same maximal value
    \item As in static optimization problems, non-uniqueness because of the
    lack of strict concavity of objective function
    \item When we include a concavity assumption, uniqueness of optimal plan
    is guaranteed
\end{itemize}
\end{frame}

\begin{frame}{Concavity Assumption}
\begin{block}{Assumption 6.3}
    $U$ is concave; that is, for any $\alpha \in (0,1)$ and any $(x,y),
    (x',y') \in X_G$:
    \begin{equation*}
        U(\alpha x + (1-\alpha)x', \alpha y + (1-\alpha)y') \geq \alpha
        U(x,y) + (1-\alpha)U(x',y')
    \end{equation*}
    and the inequality is strict if $x \neq x'$.

    In addition, $G$ is convex in the sense that for any $\alpha \in [0,1]$,
    and $x, x', y, y' \in X$ such that $y \in G(x)$ and $y' \in G(x')$:
    \begin{equation*}
        \alpha y + (1-\alpha)y' \in G(\alpha x + (1-\alpha)x').
    \end{equation*}
\end{block}
\end{frame}

\begin{frame}{Concavity of Value Function}
\begin{block}{Theorem 6.4 (Concavity of the Value Function)}
Suppose that Assumptions 6.1, 6.2, and 6.3 hold. Then the unique $V: X \to \R$ that satisfies the Bellman equation is strictly concave.
\end{block}

\begin{itemize}
    \item If Assumption 6.3 relaxed so that $U$ is concave (without
    additional requirement for $x \neq x'$), then a weaker version applies and
    implies $V$ is concave
    \item Important for uniqueness of optimal plans
\end{itemize}

\end{frame}

\begin{frame}
    \frametitle{Unique Optimal Plan and Policy Function}
\begin{block}{Corollary 6.1}
Suppose Assumptions 6.1, 6.2, and 6.3 hold. Then there exists a unique optimal plan $x^* \in \Phi(x_0)$ for any $x_0 \in X$. Moreover, the optimal plan can be expressed as:
\begin{align*}
x^*_{t+1} = \pi(x^*_t)
\end{align*}
where $\pi: X \to X$ is a continuous policy function.
\end{block}

\begin{itemize}
    \item The policy correspondence becomes a function, not just a
    correspondence.
    \item The policy function is continuous.
\end{itemize}

\end{frame}

\begin{frame}{Monotonicity Assumption}
\begin{block}{Assumption 6.4}
    For each $y \in X$, $U(\cdot, y)$ is strictly increasing in each of its
    first $K$ arguments, and $G$ is monotone in the sense that $x \leq x'$
    implies $G(x) \subset G(x')$.
\end{block}

\begin{itemize}
\item Ensures payoff function increasing in state variables
\item Larger values of state variables attractive from viewpoint of relaxing constraints
\item Natural in economic applications where ``more is better''
\end{itemize}
\end{frame}

\begin{frame}
\frametitle{Monotonicity of Value Function}

\begin{block}{Theorem 6.5 (Monotonicity)}
Suppose Assumptions 6.1, 6.2, and 6.4 hold, and let $V: X \to \mathbb{R}$ be the unique solution to the Bellman equation. Then $V$ is strictly increasing in all of its arguments.
\end{block}

\begin{itemize}
    \item Consistent with monotonic preferences
    \item Useful for comparative statics analysis
    \item In growth models, higher capital stock always increases value
    function
\end{itemize}

\end{frame}

\begin{frame}{Differentiability Assumption}
\begin{block}{Assumption 6.5}
$U$ is continuously differentiable on the interior of its domain $X_G$.
\end{block}

\begin{itemize}
    \item Common in most economic models
    \item Enables us to work with first-order necessary conditions
\end{itemize}
\end{frame}

\begin{frame}
\frametitle{Differentiability and Envelope Theorem}

\begin{block}{Theorem 6.6 (Differentiability)}
    Suppose Assumptions 6.1, 6.2, 6.3, and 6.5 hold. Let $\pi(\cdot)$ be the
    policy function from Corollary 6.1, and assume $x \in \text{Int }X$ and
    $\pi(x) \in \text{Int }G(x)$. Then $V(\cdot)$ is differentiable at $x$,
    with gradient:
    \begin{align*}
        DV(x) = D_x U(x, \pi(x))
    \end{align*}
\end{block}

\begin{itemize}
    \item Allows us to use differential calculus on the Bellman equation
    \item Enables derivation of Euler equations
\end{itemize}

\end{frame}

\begin{frame}{Summary}
    Given these assumptions, the following results are established:
    \begin{enumerate}
        \item \textbf{Theorem 6.1:} Equivalence of Values
        \item \textbf{Theorem 6.2:} Principle of Optimality  
        \item \textbf{Theorem 6.3:} Existence of Solutions
        \item \textbf{Theorem 6.4:} Concavity of Value Function
        \item \textbf{Theorem 6.5:} Monotonicity of Value Function
        \item \textbf{Theorem 6.6:} Differentiability of Value Function and
        the Envelope Theorem
    \end{enumerate}
\end{frame}

\section{The Contraction Mapping Theorem}

\begin{frame}
\frametitle{Definition}

\textbf{Metric Space:} $(S,d)$ where $S$ is a space and $d$ is a distance metric

\textbf{Operators:} Mappings $T: S \to S$ 

\begin{block}{Definition 6.1 (Contraction Mapping)}
Let $(S,d)$ be a metric space and $T: S \to S$ be an operator. If for some $\beta \in (0,1)$:
\begin{align*}
d(Tz_1, Tz_2) \leq \beta d(z_1, z_2) \text{ for all } z_1, z_2 \in S
\end{align*}
then $T$ is a contraction mapping (with modulus $\beta$).
\end{block}

\vspace{0.3cm}

\textbf{Intuition:} Contraction mappings bring elements uniformly closer together.

\end{frame}

\begin{frame}{Contraction: Example}
    
\textbf{Example 6.2:} On interval $S = [a,b]$ with usual metric $d(z_1,z_2) = |z_1 - z_2|$:
$T$ is a contraction if:
\begin{align*}
\frac{|Tz_1 - Tz_2|}{|z_1 - z_2|} \leq \beta < 1 \text{ for all } z_1 \neq z_2
\end{align*}

This means $T$ has slope less than 1 everywhere.

\end{frame}

\begin{frame}
\frametitle{The Contraction Mapping Theorem}

\begin{block}{Definition 6.2 (Fixed Point)}
A fixed point of $T$ is any element $\hat{z} \in S$ satisfying $T\hat{z} = \hat{z}$.
\end{block}


\begin{block}{Theorem 6.7 (Contraction Mapping Theorem)}
    Let $(S,d)$ be a complete metric space and suppose $T: S \to S$ is a
    contraction. Then $T$ has a unique fixed point $\hat{z}$; that is, there
    exists a unique $\hat{z} \in S$ such that:
    \begin{align*}
        T\hat{z} = \hat{z}.
    \end{align*}
    Moreover, $T^nz \to \hat{z}$ as $n \to \infty$ for any $z \in S$.
\end{block}

\begin{itemize}
    \item Simple conditions that guarantee existence, uniqueness, and convergence
    \item Applies to infinite-dimensional function spaces
    \item Foundation of dynamic programming theory
\end{itemize}

\end{frame}

% \begin{frame}
% \frametitle{Proof Idea: Contraction Mapping Theorem}

% Start with any $z_0 \in S$, define sequence:
% \begin{align*}
% z_{n+1} = Tz_n \text{ so that } z_n = T^n z_0
% \end{align*}

% \begin{enumerate}
%     \item Show $d(z_{n+1}, z_n) \leq \beta^n d(z_1, z_0)$
%     \item Prove $\{z_n\}$ is Cauchy: $d(z_m, z_n) \leq \frac{\beta^n}{1-\beta} d(z_1, z_0)$
%     \item Use completeness: $z_n \to \hat{z} \in S$  
%     \item Show $\hat{z}$ is fixed point: $T\hat{z} = \hat{z}$
%     \item Prove uniqueness by contradiction
% \end{enumerate}

% \end{frame}

\begin{frame}
\frametitle{Applications of Contraction Mappings}

\begin{block}{Theorem 6.8}
Let $(S,d)$ be complete and $T: S \to S$ be a contraction with fixed point $\hat{z}$.
\begin{enumerate}
    \item If $S'$ is a closed subset of $S$ and $T(S') \subseteq S'$, then
    $\hat{z} \in S'$
    \item If $T(S') \subseteq S'' \subseteq S'$, then $\hat{z} \in S''$
\end{enumerate}
\end{block}

\begin{itemize}
    \item If we start in a closed set that maps to itself, fixed point is in
    that set
    \item Enables proving properties like strict concavity or monotonicity
\end{itemize}

\end{frame}

\begin{frame}
\frametitle{Blackwell's Sufficient Conditions}

Checking contraction property directly is often difficult. Blackwell provides easier conditions:

\begin{block}{Theorem 6.9 (Blackwell's Conditions)}
    Let $X \subseteq \mathbb{R}^K$ and $B(X)$ be the space of bounded
    functions $f: X \to \mathbb{R}$ equipped with the sup norm. If $T: B(X)
    \to B(X)$ satisfies:
\begin{enumerate}
    \item Monotonicity: For any $f, g \in B(X)$, $f(x) \leq g(x)$ for all $x$
    implies $(Tf)(x) \leq (Tg)(x)$ for all $x \in X$; and
    \item Discounting: $\exists \beta \in (0,1)$ such that
    \begin{align*}
        [T(f + c)](x) \leq (Tf)(x) + \beta c
    \end{align*}
    for all $f \in B(X)$, $c \geq 0$, and $x \in X$.
\end{enumerate}
Then $T$ is a contraction with modulus $\beta$ on $B(X)$.
\end{block}

\end{frame}

\begin{frame}
\frametitle{Applying Blackwell's Conditions to Dynamic Programming}

For the Bellman operator:
\begin{align*}
TV(x) = \max_{y \in G(x)} \{U(x,y) + \beta V(y)\}
\end{align*}

\begin{itemize}
    \item Checking Monotonicity: If $V_1(x) \leq V_2(x)$ for all $x$, then:
    \begin{align*}
        TV_1(x) = \max_{y \in G(x)} \{U(x,y) + \beta V_1(y)\} \leq \max_{y
          \in G(x)} \{U(x,y) + \beta V_2(y)\} = TV_2(x) 
    \end{align*}
    \item Checking Discounting:
    \begin{align*}
        T(V + c)(x) &= \max_{y \in G(x)} \{U(x,y) + \beta(V(y) + c)\}\\
        &= \max_{y \in G(x)} \{U(x,y) + \beta V(y)\} + \beta c\\
        &= TV(x) + \beta c
    \end{align*}
\end{itemize}

\end{frame}

\section{Proofs of Main Theorems}

\begin{frame}
\frametitle{Key Lemma: Separating Returns}

For a feasible infinite sequence $x = (x_0, x_{1}, \ldots) \in \Phi(x_0)$,
define
\begin{equation*}
    \bar{U}(x) \equiv \sum_{t=0}^{\infty} \beta^t U(x_t, x_{t+1})
\end{equation*}

\begin{block}{Lemma 6.1}
Suppose Assumption 6.1 holds. For any $x_0 \in X$ and $x \in \Phi(x_0)$:
\begin{align*}
\bar{U}(x) = U(x_0, x_{1}) + \beta \bar{U}(x')
\end{align*}
where $x' = (x_{1}, x_{2}, \ldots)$.
\end{block}

\end{frame}

\begin{frame}
\frametitle{Proof Sketch: Theorem 6.1 (Equivalence)}

Goal: Show $V^*(x) = V(x)$ where:
\begin{itemize}
    \item $V^*(x) = \sup_{\mathbf{z} \in \Phi(x)} \bar{U}(\mathbf{z})$
    (Problem 6.2)
    \item $V(x) = \sup_{y \in G(x)} \{U(x,y) + \beta V(y)\}$ (Problem 6.3)
\end{itemize}

Idea: Use Lemma 6.1 and the definition of supremum. (Read pp. 195--197 of the
textbook.)

\end{frame}

\begin{frame}
\frametitle{Proof Sketch: Theorem 6.2 (Principle of Optimality)}

\begin{itemize}
    \item Let $x^*$ attain $V^*(x_0)$ in Problem 6.2. Our goal is to show
    that $$V^*(x^*_t) = U(x^*_t, x^*_{t+1}) + \beta V^*(x^*_{t+1})$$

    Proof idea: show that $(x^*_t, x^*_{t+1}, \ldots)$ is optimal starting
    from $x^*_t$ using Lemma 6.1 and induction
    \item Suppose $x^*$ satisfies the Bellman equation. We want to show that
    it is optimal.

    Proof idea: iterating on the Bellman equation.
\end{itemize}
\end{frame}

\begin{frame}
\frametitle{Two Proofs of Theorem 6.3 (Existence)}

\textbf{Version 1 (Sequence problem under compact $X$):}
\begin{itemize}
    \item Basic idea: show that the objective function is continuous on a
    compact set, and therefore attains a maximum
    \item Treat the objective function as continuous on $X^\infty$
    \item Show that the constraint set is a compact subset of $X^\infty$
\end{itemize}


\textbf{Version 2 (Contraction Mapping):}
\begin{itemize}
    \item Define Bellman operator $T$ on space $C(X)$ of continuous functions
    \item Show $T$ satisfies Blackwell's conditions
    \item Apply Contraction Mapping Theorem: unique fixed point exists
    \item Fixed point is solution to recursive problem
\end{itemize}

\end{frame}

\section{Applications of Stationary Dynamic Programming}

\begin{frame}{Applications}
    \begin{itemize}
        \item We show how dynamic programming can be applied to a range of
        problems
        \item \textbf{Main Result}: Theorem 6.10 shows how dynamic
        first-order conditions (Euler equations) together with the
        transversality condition are sufficient to characterize solutions
    \end{itemize}
\end{frame}

\subsection{Euler Equation}

\begin{frame}{The Functional Equation}
Consider the functional equation corresponding to Problem 6.3:
\begin{equation*}
V(x) = \max_{y \in G(x)} \{U(x,y) + \beta V(y)\} \quad \text{for all } x \in X
\end{equation*}

Throughout, we assume Assumptions 6.1–6.5 hold.
\begin{itemize}
    \item From Theorem 6.4: maximization problem is strictly concave
    \item From Theorem 6.6: maximand is differentiable
    \item For interior solutions $y \in \text{Int } G(x)$: first-order
    conditions are necessary and sufficient
\end{itemize}
\end{frame}

\begin{frame}{The Euler Equations}
For optimal solutions, we can characterize them by the \textbf{Euler equations}:
\begin{equation*}
D_y U(x, y^*) + \beta DV(y^*) = 0
\end{equation*}

\begin{itemize}
    \item Use the Envelope Theorem for dynamic programming by differentiating
    the functional equation with respect to state vector $x$:
    \begin{equation*}
        DV(x) = D_x U(x, y^*)
    \end{equation*}
    \item Substituting it into the Euler equation gives
    \begin{equation*}
        D_y U(x, \pi(x)) + \beta D_x U(\pi(x), \pi(\pi(x))) = 0
    \end{equation*}
\end{itemize}
This is a \textbf{functional equation} in the unknown function $\pi(.)$ and
characterizes the optimal policy function.
\end{frame}


\begin{frame}{One-Dimensional Case}
    When both $x$ and $y$ are real numbers, the Euler equation becomes:
    \begin{equation*}
        \frac{\partial U(x, y^*)}{\partial y} + \beta V'(y^*) = 0
    \end{equation*}
    Sum of marginal gain today from increasing $y$ plus discounted marginal
    gain from increasing $y$ on future returns must equal zero
\end{frame}

\begin{frame}{Envelope Condition - One Dimensional}
The one-dimensional Envelope Condition gives:
\begin{equation*}
V'(x) = \frac{\partial U(x, y^*)}{\partial x}
\end{equation*}

Combining with the Euler equation:
\begin{equation*}
\frac{\partial U(x, \pi(x))}{\partial y} + \beta \frac{\partial U(\pi(x), \pi(\pi(x)))}{\partial x} = 0
\end{equation*}

With time arguments explicitly:
\begin{equation*}
    \frac{\partial U(x_t, x^*_{t+1})}{\partial y} + \beta \frac{\partial
      U(x^*_{t+1}, x^*_{t+2})}{\partial x} = 0
\end{equation*}
\end{frame}

\begin{frame}{The Transversality Condition}
    \begin{itemize}
        \item Euler equation is \textbf{not sufficient} for optimality. We
        also need the transversality condition.
        \item The transversality condition ensures that there are no
        beneficial simultaneous changes in an infinite number of choice
        variables.
        \item \textbf{General case}:
        \begin{equation*}
            \lim_{t \to \infty} \beta^t D_x U(x^*_t, x^*_{t+1}) \cdot x^*_t = 0
        \end{equation*}
        \item \textbf{One-dimensional case}:
        \begin{equation*}
            \lim_{t \to \infty} \beta^t \frac{\partial U(x^*_t,
              x^*_{t+1})}{\partial x} \cdot x^*_t = 0 
        \end{equation*}
        \item The transversality condition requires that the marginal return
        from state variable $x$ times the value of this state variable does
        not increase asymptotically at rate faster than or equal to
        $1/\beta$.
    \end{itemize}
\end{frame}

\begin{frame}{Main Theorem}
\begin{block}{Theorem 6.10 (Euler Equations and Transversality Condition)}
    Let $X \subset \mathbb{R}^K_+$ and suppose Assumptions 6.1–6.5 hold. Then
    a sequence $\{x^*_t\}_{t=0}^{\infty}$ such that $x^*_{t+1} \in \text{Int
    } G(x^*_t)$, $t = 0,1,\ldots$, is optimal for Problem 6.2 given $x_0$
    if and only if it satisfies:
    \begin{align*}
        &D_y U(x^*_t, x^*_{t+1}) + \beta D_x U(x^*_{t+1}, x^*_{t+2}) = 0 \\
        &\lim_{t \to \infty} \beta^t D_x U(x^*_t, x^*_{t+1}) \cdot x^*_t = 0
    \end{align*}
\end{block}
The Euler equations together with the transversality condition are
\textbf{necessary and sufficient}, under those assumptions.
\end{frame}

\begin{frame}{Example: Optimal Growth}
    \textbf{Example 6.4}: Consider an optimal growth model with log
    preferences, Cobb-Douglas technology, and full depreciation:
    \begin{align*}
        \max_{\{k_t,c_t\}_{t=0}^{\infty}} &\sum_{t=0}^{\infty} \beta^t \log c_t \\
        \text{s.t. } &k_{t+1} = k_t^{\alpha} - c_t \\
        &k_0 > 0
    \end{align*}
    where $\beta \in (0,1)$, $k$ is capital-labor ratio, and the resource
    constraint follows from the production function $K^{\alpha}L^{1-\alpha}$
    in per capita terms.
\end{frame}

\begin{frame}{Recursive Formulation}
Write down the Bellman equation
\begin{equation*}
V(x) = \max_{y \geq 0} \{\log(x^{\alpha} - y) + \beta V(y)\}
\end{equation*}
where $x$ is today's capital stock and $y$ is tomorrow's capital stock.

\begin{itemize}
    \item Previously, we used guess and verify to solve the value function.
    \item Show that the problem satisfies Assumptions 6.1--6.5, so Theorems
    6.1--6.6 apply. Then we can also use the Euler equation.
\end{itemize}
\end{frame}

\begin{frame}{Applying Euler Equations}
    \begin{itemize}
        \item Since $V(.)$ is differentiable, the first-order condition is
        \begin{equation*}
            \frac{1}{x^{\alpha} - y} = \beta V'(y)
        \end{equation*}
        \item Envelope Condition:
        \begin{equation*}
            V'(x) = \frac{\alpha x^{\alpha-1}}{x^{\alpha} - y}
        \end{equation*}
        \item Combining using $y = \pi(x)$:
        \begin{equation*}
            \frac{1}{x^{\alpha} - \pi(x)} = \beta \frac{\alpha
              \pi(x)^{\alpha-1}}{\pi(x)^{\alpha} - \pi(\pi(x))}
        \end{equation*}
    \end{itemize}
    This is a functional equation in a single function $\pi(x)$.
\end{frame}

\begin{frame}{Guess-and-Verify Solution}
    \begin{itemize}
        \item There are no straightforward ways to solve functional
        equations, but guess-and-verify often works.
        \item Let us conjecture that $\pi(x) = ax^{\alpha}$
        \item Substituting into the Euler equation:
        \begin{align*}
            \frac{1}{x^{\alpha} - ax^{\alpha}} &= \beta \frac{\alpha
              a^{\alpha-1} x^{\alpha(\alpha-1)}}{a^{\alpha} x^{\alpha^2} -
              a^{1+\alpha} x^{\alpha^2}} = \frac{\beta}{a}
            \frac{\alpha}{x^{\alpha} - ax^{\alpha}}
        \end{align*}
        which implies $a = \beta\alpha$ satisfies the equation.
    \end{itemize}
\end{frame}

\begin{frame}{Complete Solution}

\begin{itemize}
    \item The policy function is: $\pi(x) = \beta\alpha x^{\alpha}$
    \item Law of motion: $k_{t+1} = \beta\alpha k_t^{\alpha}$
    \item Optimal consumption: $c_t = (1-\beta\alpha)k_t^{\alpha}$
    \item Capital-labor ratio $k_t$ converges to steady state $k^*$, which
    ensures transversality condition
    \item By Corollary 6.1 and Theorem 6.10: $\pi(x) = \beta\alpha
    x^{\alpha}$ is the unique policy function
\end{itemize}
\end{frame}

\begin{frame}{Example: Optimal Savings Problem}
    \textbf{Example 6.5}: Consider an infinitely-lived consumer who solves
    \begin{align*}
        &\max_{\{c_t,a_t\}_{t=0}^{\infty}} \sum_{t=0}^{\infty} \beta^t u(c_t) \\
        &\text{subject to: } a_{t+1} = (1+r)a_t + w - c_t, \quad c_t \geq 0
    \end{align*}
    \begin{itemize}
        \item $u(c)$ is strictly increasing, continuously differentiable,
        strictly concave
        \item Without further constraints, this is not well-defined and
        allows consumer to build unlimited debt ($\lim_{t \to \infty} a_t =
        -\infty$)---``Ponzi games''.
    \end{itemize}
\end{frame}

\begin{frame}{How to Avoid Ponzi Games}
Three approaches to prevent unlimited borrowing:
\begin{enumerate}
\item \textbf{No-Ponzi condition}: Rule out such schemes directly
\item \textbf{No borrowing}: $a_{t+1} \geq 0$ for all $t$ (too restrictive)
\item \textbf{Natural debt limit}: Maximum repayable debt $a_{t+1} \geq
\underline{a} \equiv -w/r$
\end{enumerate}
\end{frame}

\begin{frame}{Ensuring Compactness}
    \begin{itemize}
        \item \textbf{Challenge}: Assets $a$ don't necessarily belong to
        a compact set.
        \item \textbf{Solution}: Choose upper bound $\bar{a}$ and
        restrict $a \in [\underline{a}, \bar{a}]$.
        \begin{enumerate}
            \item Solve problem on a compact set
            \item Verify that the solution indeed lies in the state space
        \end{enumerate}
        \item \textbf{Natural choice}: $\bar{a} \equiv a_0 + w/r <
        \infty$
    \end{itemize}
\end{frame}

\begin{frame}{Recursive Formulation}
    \begin{itemize}
        \item State variable: $a_t$
        \item Consumption: $c_t = (1+r)a_t + w - a_{t+1} \geq 0$
        \item The Bellman equation
        \begin{equation*}
            V(a) = \max_{a' \in [\underline{a}, \bar{a}],\, a' \leq (1+r)a+w}
            \{u((1+r)a + w - a') + \beta V(a')\}
        \end{equation*}
    \end{itemize}
    Verify that the assumptions are satisfied.
\end{frame}

\begin{frame}{Consumption Euler Equation}
    \begin{itemize}
        \item Using one-dimensional Euler equation and the Envelope Condition gives
        \begin{equation*}
            u'(c) = \beta(1+r)u'(c')
        \end{equation*}
        \item Since $u'(.)$ exists and is strictly decreasing (because $u$ is
        continuously differentiable and strictly concave), we get the
        following simple rule:
        \begin{align*}
            \text{if } r &= \beta^{-1} - 1, c = c' \text{ (constant consumption)} \\
            \text{if } r &> \beta^{-1} - 1, c < c' \text{ (increasing consumption)} \\
            \text{if } r &< \beta^{-1} - 1, c > c' \text{ (decreasing consumption)}
        \end{align*}
    \end{itemize}
\end{frame}

\subsection{Dynamic Programming vs Sequence Problem}

\begin{frame}{Transversality Condition from Finite Horizon Problem}
    Consider a dynamic optimization problem with finite horizon $T$:
    \begin{align*}
        &\max_{\{x_t\}_{t=1}^{T+1}} \sum_{t=0}^T \beta^t U(x_t, x_{t+1}) \\
        &\text{s.t. } x_{t+1} \geq 0, \text{ } x_0 \text{ given}
    \end{align*}
    \vspace{-2ex}
    \begin{itemize}
        \item When $0 \leq t \leq T-1$, FOCs from finite-dimensional
        optimization assuming interior solution:
        \begin{equation*}
            \frac{\partial U(x^*_t, x^*_{t+1})}{\partial y} + \beta
            \frac{\partial U(x^*_{t+1}, x^*_{t+2})}{\partial x} = 0 
        \end{equation*}
        
        \item When $t = T$, the complementary slackness condition implies:
        \begin{equation*}
            x^*_{T+1} \geq 0, \text{ and } \beta^T \frac{\partial U(x^*_T,
              x^*_{T+1})}{\partial y} x^*_{T+1} = 0 
        \end{equation*}
    \end{itemize}
\end{frame}

\begin{frame}{Example: Finite Horizon Optimal Growth}
    \textbf{Example 6.6:} In the optimal growth problem,
    \begin{equation*}
        U(x_t, x_{t+1}) = u(f(x_t) + (1-\delta)x_t - x_{t+1})
    \end{equation*}
    Suppose that the world ends at date $T$. Then at the last date $T$:
    $$\frac{\partial U(x^*_T, x^*_{T+1})}{\partial y} = -u'(c^*_T) < 0$$
    From the boundary condition, an optimal path must have $k^*_{T+1} = x^*_{T+1} = 0$.

    No capital should be left at end of world. Any leftover resources should
    be consumed.
\end{frame}

\begin{frame}{Deriving Transversality Condition}
    \begin{itemize}
        \item The transversality condition can be derived heuristically by
        taking the limit of the boundary condition
    \begin{equation*}
        \lim_{T \to \infty} \beta^T \frac{\partial U(x^*_T,
          x^*_{T+1})}{\partial y} x^*_{T+1} = 0
    \end{equation*}
    \item The Euler equation implies
    \begin{equation*}
        \frac{\partial U(x^*_T, x^*_{T+1})}{\partial y} + \beta
        \frac{\partial U(x^*_{T+1}, x^*_{T+2})}{\partial x} = 0
    \end{equation*}
    \item Substituting and changing timing:
    \begin{equation*}
        \lim_{T \to \infty} \beta^T \frac{\partial U(x^*_T, x^*_{T+1})}{\partial x} x^*_T = 0
    \end{equation*}
    which is exactly the transversality condition: a ``boundary condition at
    infinity.''
    \end{itemize}
\end{frame}

\end{document}
